
\section{Introduction}

The architecture of global financial intermediation is undergoing a fundamental structural transition. Traditional centralized financial (TradFi) systems, characterized by institutional intermediaries and lagged regulatory reporting, are increasingly complemented or replaced by Decentralized Finance (DeFi) protocols \cite{}. This emerging ecosystem utilizes permissionless blockchains and autonomous smart contracts to replicate core financial functions, including lending, trading, and asset management \cite{}. However, the rapid expansion of these networks has created a critical measurement gap. While blockchain data is inherently transparent and real-time, the complex, multi-layered composability of DeFi protocols obscures the true nature of inter-protocol credit exposures and systemic dependencies \cite{[8]}.

In this paper, we introduce \textbf{DeXposure-FM}, a multimodal foundation model designed to measure and forecast credit exposures across the global DeFi landscape. Unlike traditional econometric models that often treat financial sectors in isolation, DeXposure-FM treats the entire ecosystem as a unified "world model" \cite{}. The model leverages the \textbf{DeXposure} dataset, a large-scale reconstruction of value-linked credit exposures encompassing 43.7 million entries across 4,300+ protocols on 602 blockchains, covering the period from 2020 to 2025 \cite{[15, 17, 19]}. 

The core technical contribution of DeXposure-FM lies in its joint architectural design, which simultaneously processes time-series dynamics, graph topologies, and tabular protocol metadata \cite{}. By framing financial measurement as a foundation model task, we enable high-fidelity "zero-shot" forecasting of protocol-level Total Value Locked (TVL), the weights of credit-exposure links, and structural shifts in sector-level connectivity \cite{}. To ensure the model generalizes across non-stationary market regimes, we employ \textbf{Sharpness-Aware Minimization (SAM)} \cite{}. This optimizer seeks parameters in flatter regions of the loss landscape, providing robust performance even during extreme stress events such as the collapse of Terra/Luna or the failure of major exchange platforms \cite{}.

Beyond machine learning benchmarks, we utilize DeXposure-FM to advance the frontier of macroprudential monitoring. We construct dynamic measures of protocol-level systemic importance, sector-to-sector spillover indices based on the Diebold-Yilmaz framework, and early-warning indicators derived from shifts in network concentration and collateral fragility \cite{[5, 16, 18, 37]}. Our results demonstrate how AI-generated statistics can support counterfactual policy evaluations—such as the impact of caps on leverage or bridge volumes—while providing a rigorous basis for central banks and regulators to interpret the systemic risks inherent in digital financial infrastructure \cite{[7, 40, 49]}.

